\documentclass[11pt,a4paper]{article}
\usepackage[utf8]{inputenc}
\usepackage[T1]{fontenc}
\usepackage{amsfonts}
\usepackage[french]{babel}
\frenchbsetup{StandardLists=true} % à inclure si on utilise \usepackage[francais]{babel}
\usepackage{enumitem}
\usepackage{amssymb}
\usepackage{amsmath}
\usepackage[left=2cm,right=2cm,top=2cm,bottom=2cm]{geometry}
\usepackage{multirow}
\usepackage[dvipsnames]{xcolor}

\usepackage[T1]{fontenc}
\usepackage{fouriernc}
\usepackage{graphicx}

\usepackage{setspace}
\singlespacing

\usepackage{pstricks}
\usepackage{xkeyval}
\usepackage{pst-labo}


\usepackage{multicol}
\usepackage{caption}
\usepackage{fancyhdr}
\pagestyle{fancy}
\renewcommand\headrulewidth{1pt}
\fancyhead[L]{Lycée Denis Diderot }
\fancyhead[R]{Classe de TSN1}
\setcounter{page}{3}\fancyfoot[C]{page \thepage}
\fancyfoot[R]{Année 2025 2026}
\fancyfoot[L]{Alexis RUIZ}
\usepackage{multirow,tabularx,booktabs}
\newcommand{\cadreblanc}[1]{%
\medskip\framebox[\linewidth][c]{%
\rule{0mm}{#1mm}}\par
\medskip}

\usepackage{awesomebox}
\setlength{\aweboxvskip}{0mm}
\usepackage{pifont}
\usepackage{chemmacros}
\chemsetup{modules=all}
\setlength{\columnseprule}{.25mm}



\newcommand*\acocher{\quitvmode{\fboxsep0pt \fboxrule0.8pt \fbox{\vrule height1.8ex depth.1ex width0pt \vrule height0pt depth0pt width1.9ex }}}
\newcolumntype{R}[1]{>{\raggedleft\arraybackslash }b{#1}}
\newcolumntype{L}[1]{>{\raggedright\arraybackslash }b{#1}}
\newcolumntype{C}[1]{>{\centering\arraybackslash }b{#1}}

\newcommand{\defproblem}{\awesomebox[black]{2pt}{\rotatebox{90}{\large{\textbf{Problématique}}}}{black}}
\newcommand{\defconclusion}{\awesomebox[black]{2pt}{\rotatebox{90}{\Large{\textbf{Conclusion}}}}{black}}

\frenchbsetup{StandardLists=true}
\begin{document}



\begin{center}
\LARGE{EXERCICES}
\end{center}




\hrule

\vspace{1cm}

\textbf{Exercice 1 - QCM }
Conversions d'unités de longueurs multiples et sous-multiples :

\begin {multicols}{3}

\begin{enumerate}
\item Combien de millimètres y a-t-il dans un mètre ?
\begin{enumerate}
\item[a.] 10
\item[b.] 100
\item[c.] 1000
\item[d.] 10000
\end{enumerate}
\item Combien de micromètres y a-t-il dans un mètre ?
\begin{enumerate}
\item[a.] 10
\item[b.] $10^2$
\item[c.] $10^4$
\item[d.] 1 000 000
\end{enumerate}
\item Combien de kilomètres y a-t-il dans un millimètre ?
\begin{enumerate}
\item[a.] $10^{-6}$
\item[b.] $10^{-5}$
\item[c.] $10^{-4}$
\item[d.] $10^{-3}$
\end{enumerate}

\columnbreak

\item Combien de nanomètres y a-t-il dans un millimètre ?
\begin{enumerate}
\item[a.] $10^3$
\item[b.] $10^6$
\item[c.] $10^9$
\item[d.] $10^{12}$
\end{enumerate}
\item Combien de centimètres y a-t-il dans un kilomètre ?
\begin{enumerate}
\item[a.] 100
\item[b.] 1000
\item[c.] 10000
\item[d.] 100000
\end{enumerate}
\item Combien de picomètres y a-t-il dans un centimètre ?
\begin{enumerate}
\item[a.] $10^{10}$
\item[b.] $10^{12}$
\item[c.] $10^{14}$
\item[d.] $10^{16}$
\end{enumerate}
\item Combien de décimètres y a-t-il dans un kilomètre ?
\begin{enumerate}
\item[a.] 100
\item[b.] 1000
\item[c.] 10000
\item[d.] 100000
\end{enumerate}

\columnbreak


\item Combien de femtomètres y a-t-il dans un mètre ?
\begin{enumerate}
\item[a.] $10^{12}$
\item[b.] $10^{13}$
\item[c.] $10^{14}$
\item[d.] $10^{15}$
\end{enumerate}
\item Combien de décamètres y a-t-il dans un hectomètre ?
\begin{enumerate}
\item[a.] 1
\item[b.] 10
\item[c.] 100
\item[d.] 1000
\end{enumerate}
\item Combien de attomètres y a-t-il dans un millimètre ?
\begin{enumerate}
\item[a.] $10^{12}$
\item[b.] $10^{15}$
\item[c.] $10^{18}$
\item[d.] $10^{21}$
\end{enumerate}
\end{enumerate}

\end{multicols}


\dotfill \textbf{Notes  } \dotfill

\begin{multicols}{2}

1 millimètre (mm) = $10^{-3}$ mètres (m)

1 micromètre ($\mu$m) = $10^{-6}$ mètres (m)

1 nanomètre (nm) = $10^{-9}$ mètres (m)

1 picomètre (pm) = $10^{-12}$ mètres (m)

1 femtomètre (fm) = $10^{-15}$ mètres (m)

1 attomètre (am) = $10^{-18}$ mètres (m)

1 kilomètre (km) = 1000 mètres (m)

1 centimètre (cm) = 0.01 mètres (m)

1 décimètre (dm) = 0.1 mètres (m)

1 décamètre (dam) = 10 mètres (m)

1 hectomètre (hm) = 100 mètres (m)

\end{multicols}

\newpage

\newpage

\textbf{Exercice 2 - QCM}\\
Cocher la ou les bonnes réponses\\
\begin{multicols}{2}
 La célérité du son dans l'air à $20^{\circ} C$ est : 
\begin{itemize}[label=$\square$]
	\item 343 km/s
	\item 343 m/s
	\item 343 m/h
\end{itemize}


\ding{203} La célérité de la lumière est de 300 000 km/s, donc la lumière est : 
\begin{itemize}[label=$\square$]
	\item plus lente que le son
	\item de même rapidité que le son
	\item plus rapide que le son
\end{itemize}


\ding{204} Le son est la lumière peuvent se déplacer dans le vide : 
\begin{itemize}[label=$\square$]
	\item vrai
	\item faux
	\item vrai pour la lumière seulement
\end{itemize}


\ding{205} La La relation faisant le lien entre la période $T$, la longueur d'onde $\lambda$ et la vitesse du son $v$ est : 
\begin{itemize}[label=$\square$]
	\item $\lambda = v \times T$
	\item $v = \lambda \times T$
	\item $T = \lambda \times v$
\end{itemize}

\columnbreak

\ding{206} La longueur d'onde s'exprime en  : 
\begin{itemize}[label=$\square$]
	\item secondes (s)
	\item mètres (m)
	\item hertz (Hz)
\end{itemize}

\ding{207} L'oreille possède une membrane qui vibre lorsqu'elle reçoit un son. Cette membrane s'appelle : 
\begin{itemize}[label=$\square$]
	\item le pavillon
	\item le tympan
	\item le conduit auditif
\end{itemize}

\ding{208} L'exposition de l'oreille à un son dont le niveau sonore est supérieur à 120 dB est : 
\begin{itemize}[label=$\square$]
	\item sans danger
	\item dangereuse
	\item responsable de la destruction irréversible de l'audition
\end{itemize}

\ding{209} Le son est plus rapide : 
\begin{itemize}[label=$\square$]
	\item dans l'eau
	\item dans l'air
	\item dans un métal
\end{itemize}

\ding{210}  Si je double la distance entre mon oreille et la source sonore, le niveau sonore est :
\begin{itemize}[label=$\square$]
	\item diviser par deux
	\item abaissé de 6 dB
	\item augmente de 4 dB
\end{itemize}

\end{multicols}

\newpage

\textbf{Exercice 3 - Mesure de la vitesse du son}
\begin{multicols}{2}
En cours de sciences physiques Eliott utilise un haut parleur et deux microphones placées côte à côte. Il obtient deux signaux en phase sur l'oscilloscope. Puis il déplace l'un des micros d'une distance $\lambda$ afin d'obtenir à nouveau les deux signaux en phase. 
\begin{center}
    
\def\scl{0.5}%scaling factor of the picture
\begin{tikzpicture}[
  scale=\scl,
  controlpanels/.style={yellow!30!,rounded corners,draw=black,thick},
  screen/.style={white,draw=black,thick},
  trace1/.style={Blue , ultra thick},
  trace2/.style={Red , ultra thick},
  %smallbutton/.style={white,draw=black, thick},
  axes/.style={thick}]

  % Screen, centered around the origin then shifted for easy plotting
  \begin{scope}[xshift=7cm,yshift=8cm,samples=150]
    \fill[black!30!,rounded corners,draw=black,thick](-5.3,-4.3)
      rectangle (5.3,4.3);
    \fill[screen] (-5.0,-4.0) rectangle (5.0,4.0);
    
    \draw[trace1] plot(\x,{3*sin((1.57*\x) r});
    \draw[trace2] plot(\x,{2*sin((1.57*\x+2) r});
    
    \draw[thin] (-5.0,-4.0) grid (5.0,4.0);
    \draw[axes] (-5,0)--(5,0); % Time axis
    \draw[axes] (0,-4)--(0,4);
    \foreach \i in {-4.8,-4.6,...,4.8} \draw (\i,-0.1)--(\i,0.1);
    \foreach \i in {-3.8,-3.6,...,3.8} \draw (-0.1,\i)--(0.1,\i);
  \end{scope}

  
\end{tikzpicture}
\end{center} 
\end{multicols}
\ding{202} Sur l'oscilloscope, le balayage est de $5 \mu s $/ division , déterminer la période $T$ en $s$.\\
\cadreblanc{25}

\ding{203} Il faisait environ $20 ^\circ C $ dans la salle. \`A partir de la relation $\lambda = c_{son} \times T$, calculer la longueur d'onde $\lambda$  du son utilisé.\\
\cadreblanc{25} 


\textbf{Exercice 4 - Seuil de douleur}\\
L'oreille humaine est capable de détecter des sons faibles, mais aussi de détériorer si l'intensité acoustique devient trop forte. Le seul de douleur apparait pour une valeur d'intensité égale à 1W/$m^2$.\\[3pt]
En utilisant la relation $L=10 log \left(\dfrac{I}{I_0}\right)$ et $I_0=10^{-12}W/m^2$.\\

\ding{202} Calculer le niveau sonore $L_1$ lié au seuil de douleur en dB.\\
\cadreblanc{12}
\ding{203} Calculer le niveau sonore $L_2$ lié à une intensité de $I=3,16 \times 10^{-4} W/m^2$. Arrondir à l'unité. \\
\cadreblanc{12}
\ding{204} A quel niveau correspond cette valeur ?\\
\cadreblanc{15}


\end{document}